\chapter{Clause Linking and Converb Constructions}\label{ch:clause-linking}\label{ch:complex}

\section{Types of clause combining}
\label{sec:th-complex-overview}

Dependent clauses precede the material on which they depend: finite
complements precede the matrix predicate, converb adjuncts precede the matrix
clause, and relative modifiers precede the nominal head. These orders extend
Iridian's predominantly head-final organisation across clausal boundaries,
although complementation and adjunction involve different structural relations.
Table~\ref{tab:th-combining-strategies} gives the main combining strategies and
their formal properties.

\begin{table}[ht]
	\footnotesize
	\caption{Main clause-combining strategies in Iridian.}\label{tab:th-combining-strategies}
	\begin{tabularx}{\textwidth}{l L L l}
		\toprule
		{\sc strategy} & {\sc formal means} & {\sc function} & {\sc ref.} \\
		\midrule
		Converb chaining &
		Non-finite converb suffixes (\ird{-u}, \ird{-ení}, \ird{-ěc},
		\ird{-ic}, \ird{-it}, \ird{-ouc}, \ird{-ema}, etc.) &
		Adverbial subordination: temporal, causal, conditional,
		concessive, purposive, contrastive &
		§\,\ref{sec:converb-chains} \\[4pt]
		Finite \ird{po}-complementation &
		Finite embedded clause + linker \ird{-e} + \ird{po} (or
		support \ird{ce} + \ird{po} with zero-copula predicates) &
		Propositional complementation: saying, thinking, knowing,
		believing, etc. &
		§\,\ref{sec:complementation} \\[4pt]
		Non-assertive \ird{poby} &
		Finite embedded clause + linker \ird{-e} + \ird{poby} (or
		support \ird{ce} + \ird{poby} with zero-copula predicates) &
		Non-assertive finite complementation: embedded polar issues,
		embedded wh-content questions, and related indirect
		interrogative content &
		§\,\ref{sec:poby} \\[4pt]
		Purposive converb complement &
		Purposive converb + matrix-pivot control &
		Event complementation: wanting, trying, intending, beginning &
		§\,\ref{sec:purposive} \\[4pt]
		Coordination &
		\ird{a}, \ird{še}, \ird{má}/\ird{ozná}, disjunctives &
		Symmetric clause-level coordination &
		§\,\ref{sec:coordination} \\[4pt]
		Relativisation &
		Prenominal participant-oriented relative in \ird{-ový} &
		Nominal modification by a clausal relative &
		§\,\ref{sec:th-relativization} \\
		\bottomrule
	\end{tabularx}
\end{table}

\subsection{Head-finality in clause combining}
\label{sec:th-subord-headfinal}

The unmarked pattern in clause combining is as follows:

\begin{center}
\term{[subordinate clause] \quad [matrix-clause material] \quad [matrix verb]}
\end{center}

The subordinate clause occupies the left side of the construction; the matrix
predicate closes it. This ordering holds across all of the major subordinate
clause types: converb chains, finite complement clauses, embedded wh-clauses,
and relative clauses all precede the higher predicate they are embedded under.
The same head-final logic applies inside the subordinate clause itself, which is
internally verb-final.

\begin{figure}[ht]
	\centering
	\begin{forest}
		for tree={
			align=center,
			parent anchor=south,
			child anchor=north,
			l sep=12pt,
			s sep=16pt,
		}
		[S
			[{XP\\{\footnotesize subordinate clause}}]
			[\synhead{VP}
				[{XP\\{\footnotesize matrix-clause material}}]
				[{V\\{\footnotesize matrix}}]
			]
		]
	\end{forest}
	\caption{Head-final organisation of subordinate constructions.}
\end{figure}

There is a practical consequence for information structure. In a converb chain,
all aspectual contrasts and alignment contrasts are borne by the single finite
verb at the end of the chain; converbs are aspectually neutral and take
nominative-accusative alignment regardless of the aspect that the matrix verb
will bear. This creates an asymmetry between converb-internal and matrix-level
information marking that is discussed in §\,\ref{sec:converb-chains}.

\subsection{Subordination types: a typological note}
\label{sec:th-subord-types}

Following standard typological distinctions (Haspelmath 1995; Nikolaeva 2007),
Iridian clause-combining strategies may be classified along two parameters: the
\emph{morphological type of the dependent verb} (finite or non-finite) and the
\emph{semantic relationship} between dependent and matrix clause.

Non-finite (converb) subordination is the dominant strategy for adverbial
clauses. Finite subordination by means of \ird{po}/\ird{poby} is restricted to
argument clauses (complement clauses). Prenominal relatives employ the dedicated head \ird{-ový}, which selects
an aspect-bearing participant description without an overt relative pronoun
or complementiser (§\,\ref{sec:th-relativization}). Their embedded aspect
distinguishes them from the aspectually neutral adverbial converbs.

The dominance of converb-based strategies means that what in a European language
would be expressed as two (or more) finite clauses typically appears in Iridian
as a converb chain followed by a single finite verb. The aspectual and alignment
information that would be distributed across multiple finite verbs in a Slavic
language is therefore concentrated at the right edge of the chain. This property
has implications for switch-reference tracking (§\,\ref{sec:switch-reference})
and for discourse structure (Chapter~\ref{ch:discourse}).

\section{Converbs}\label{sec:converb-morphology}


Converb suffixes replace the aspect-alignment slot entirely. A verb carrying a
converb suffix accordingly bears no finite viewpoint or mood exponent and
no aspect-conditioned alignment contrast. Its clause nevertheless licenses
nominal case through the converb construction (§\,\ref{sec:vm-converb-case}). The
morphological structure of a converb is:

\medskip
\begin{center}
(\textsc{lex-prefix}) + \Stem{} + (\Appl{}) + (\Caus{}) + \textsc{converb-suffix}
\end{center}
\medskip

The converb suffixes and the stem mutations they occasion are set out in the following table:

\begin{table}[ht]
	\footnotesize
	\caption{Converb suffixes and their phonological effects.}\label{tab:converb-morph}
	\begin{tabularx}{\textwidth}{L L L L}
		\toprule
		{\sc converb} & {\sc suffix(es)} & {\sc stem mutation} & {\sc notes} \\
		\midrule
		Sequential    & \ird{-u}, \ird{-ení} & none & \ird{-ní} after applicative; full \ird{-ení} after causative \\
		Simultaneous/manner & \ird{-ěc} & palatalisation & ordinary consonant mutation; independent stem-grade selection \\
		Temporal (6)  & \ird{-ezak}, \ird{-eškady}, \ird{-eeby}, \ird{-ešhoume}, \ird{-eebymy}, \ird{-esny} & none & all add \ird{-e-} linker \\
		Concessive    & \ird{-ouc} & none & \\
		Conditional   & \ird{-ic}, \ird{-emy}, \ird{-ezmy}, \ird{-ena} & none for \ird{-ic}; see note & \ird{-ic} does not trigger palatalisation \\
		Purposive     & \ird{-it} & irregular; mainly Class I & \\
		Contrastive   & \ird{-ema} & none & \\
		\bottomrule
	\end{tabularx}
\end{table}

The simultaneous/manner converb \ird{-ěc} selects the front grade,
irrealis \ird{-ěv} the mixed grade, and antipassive \ird{-aš-} its separate
replacement classes. Conditional \ird{-ic} selects the plain grade
(§\,\ref{sec:cond-nopala} of Chapter~\ref{ch:complex}); its initial vowel
leaves alveolar sibilants unchanged. Sequential \ird{-ení} is also
plain-grade, unlike the homophonous front-coded nominal ending. The shortened
\ird{-ní} is selected immediately after applicative \ird{-éb-}; a following
rhotic causative instead selects full \ird{-ení}, as in \ird{pjáštrení}
and \ird{pjaštébrení}, beside the alternative in \ird{-u}. Simultaneous
\ird{-ěc} leaves causative /r/ unchanged and is written with plain
\ird{e}: \ird{pjáštrec}.
Grade membership must therefore be recorded with each construction.


\subsection{Formal analysis: the syntactic status of converbs}
\label{sec:vm-converbs}
\label{sec:vm-converb-aspless}\label{sec:vm-converb-case}\label{sec:vm-converb-adjunct}


Converbs lack finite Asp.Align and its viewpoint distinctions, while retaining
the lower predicate and any licensed applicative or causative structure.
The analysis treats the converb suffix as the exponent of a converb head,
Conv, with a defective aspectual specification. Its features specify the
linking relation without supplying
the finite IPFV/ERG-domain opposition. Conv selects the lower verbal
predicate directly, excluding the finite perfective/retrospective configuration
that licenses ergative alignment.
The independent absence of outer patient-selecting Voice, including both
ordinary \ird{-n-} and causative-conditioned \ird{-an-}, from the
established converb forms is described in §\,\ref{sec:converb-chains};
its exclusion does not follow from the absence of ergative case alone, since
the same Voice head also occurs in finite nominative-domain passives.


Within this reduced structure, converb clauses license the NOM-domain case frames without finite viewpoint
aspect. An ordinary active selects its agent or sole participant as direct and
assigns a retained patient indirect case; a simple applicative selects the
applied participant, with actor and retained theme indirect, while a higher
causative selects its causer. No converb licenses transitive-case agents.
These are constructional licensing conditions: absence of perfective or
retrospective Asp excludes the ERG frame, but does not by itself determine
which participant is privileged.


At the clause-combining level, converb clauses function as adverbial
adjuncts to the matrix clause. They are merged as adjuncts to the matrix
\synhead{IP} or \synhead{vP}, not as arguments in Spec or Complement positions.
This explains their preverbal (and pre-matrix) position: as adjuncts to
\synhead{IP}, they precede all material within \synhead{IP} and surface before
the matrix arguments in the linear string, consistent with the head-final
adjunction pattern.

Reference across converb chains follows a default same-pivot relation
 (§\,\ref{sec:switch-reference}). Distinct overt nominals or discourse
 context can support different-pivot readings, without a switch-reference
 clitic. Since an aboutness topic need not be the absolutive pivot, topic
 continuity and this dependency are separate relations. These constructions
 do not by themselves establish the availability of embedded hanging topics
 or high topicalisation in every non-finite clause.

\section{Converb chains: structure and alignment}\label{sec:converb-chains}

\subsection{The finite-verb-final rule}

In a converb chain, all verbs except the last are non-finite (converbs). Only
the final verb is finite and carries aspect, mood, and alignment. This rule is
absolute in standard written Iridian; in colloquial and dialectal speech, finite
verbs may occasionally appear in non-final positions in certain coordinated
structures.

\subsection{Prosodic integration of the chain}\label{sec:converb-phrasing}

A tightly integrated converb chain can form one intonational phrase containing
several phonological phrases (§\,\ref{sec:prosodic-domains}). Each lexical
word retains its own initial primary stress and subsidiary footing; integration
at phrase level does not create a single lexical stress domain for the chain.
A nonfinal phonological phrase may carry the continuation tone H-, while the
final boundary belongs to the intonational phrase as a whole
(§\,\ref{sec:intonation}). A converb ending therefore requires neither a
pause nor a full pitch reset.

More independent discourse units may form separate intonational phrases, whose
boundaries can reinforce a shift of topic or a new discourse move. Their
placement depends on information structure and speech rate as well as syntactic
organisation, so neither a clause boundary nor a change of pivot alone fixes
an intonational boundary. The independent contour of an insubordinate
conditional is a construction-specific case (§\,\ref{sec:cond-insub}).

\subsection{Accusative alignment within converb clauses}

Converb clauses belong to the nominative-accusative domain. In an ordinary
active converb the agent or sole participant is direct and a retained patient
is indirect, while a simple applicative selects its applied participant as
direct and assigns indirect case to the remaining arguments. A higher
causative selects the causer as direct. Converbs carry no finite viewpoint
aspect, so they license neither transitive-case agents nor the
ergative-domain antipassive.

\subsection{Switch-reference and topic continuity}\label{sec:switch-reference}

By default, the absolutive pivot of a converb clause is coreferent with the
absolutive pivot of the following clause. Because ordinary active converb
clauses privilege their agents, this default often resembles a same-subject
(SS) pattern, while the broader generalisation concerns \term{same-pivot}
 continuity. A discourse topic can be distinct from either pivot, so topical
 continuity does not by itself determine converb coreference. When
the pivots differ (different-subject/different-pivot, DS), the contrast is
signalled by overt lexical NPs, demonstratives, or the discourse structure
itself; Iridian has no grammaticalised switch-reference clitic. The presence of a
high non-direct topic does not independently change this default relation.
The distribution of linked topicalisation and external hanging topics inside
converb clauses is not specified by the same-pivot rule. Where both
clauses supply full nominal pivots, disambiguation is immediate; where one pivot
is omitted, recoverability from discourse context determines the SS/DS reading.

\section{Sequential constructions}\label{sec:sequential}

\subsection{Forms: \ird{-u} and \ird{-ení}}

The sequential converb has two forms with subtly different temporal force. The
\ird{-u} form implies a temporally bounded, completed sub-event preceding the
main clause event. The \ird{-ení} form is temporally underspecified and is
compatible with both sequential and simultaneous readings.

\pex
\a
\begingl
\gla Marek zdravení Janek znóvime.//
\glb Marek sleep-\Seq{} Janek study-\Prog{}//
\glft \trsl{While Marek sleeps, Janek is studying.} \lingcontext{simultaneous reading with \ird{-ení}}//
\endgl
\xe

\subsection{The copula converb and the discourse particle \ird{nu}}\label{sec:seq-copula}

The marked copular verb \ird{neh-} (see §\,\ref{sec:neh} of
Chapter~\ref{ch:predicates}) has the sequential converb form \ird{nesu}, a
suppletive form established independently of the regular aspect paradigm. In
colloquial speech, \ird{nesu} is contracted to \ird{nu}. The form \ird{nu} has
been grammaticalised in some dialects as a discourse particle inserted between
clauses in the manner of English `and then', without strong sequential
implication and without retaining the inchoative/copular force of the underlying
verb.

\subsection{Interaction with applicative and causative}

\pex
\a \ird{pjašt-u} / \ird{pjašt-ení} \trsl{having eaten / and [then] ate\ldots}
\a \ird{pjašt-éb-u} / \ird{pjašt-éb-ní} \trsl{having eaten for someone's benefit}
\a \ird{pjášt-r-u} / \ird{pjášt-r-ení} \trsl{having made someone eat}
\xe

The sequential suffix retains full \ird{-ení} after rhotic causative
morphology, yielding \ird{pjáštrení}; shortened \ird{-ní} occurs
after an immediately preceding applicative \ird{-éb-}. Thus an applied
causative has \ird{pjaštébr-ení}, with the causative determining the
outer suffixal environment (§\,\ref{sec:converb-morphology}).

\section{Temporal constructions}\label{sec:temporal}

\subsection{The simultaneous / manner converb (\ird{-ěc})}

The simultaneous converb (\ird{-ěc}) describes an action that is ongoing or
co-occurring with the main clause event, and also functions as a manner converb.
The suffix selects the palatal/front consonant grade
(§\,\ref{sec:suffix-classes}), independently of the sibilant-replacement
system of the antipassive. When the converb and main clause have different
subjects, the sequential \ird{-ení} is preferred for purely simultaneous
meaning.

This converb should not be confused with the derived manner-adverb suffix
\ird{-ě} discussed in §\,\ref{sec:stative-adv}. The converb in \ird{-ěc}
remains a non-finite verbal form used for clause-linking and event-structural
modification, whereas forms such as \ird{skorně} \trsl{quickly} are lexical
adverbs, not reduced clauses and not members of the converb paradigm.

\subsection{The six temporal converbs}

\begin{table}[ht]
	\footnotesize
	\caption{Temporal converb suffixes.}\label{tab:tempconverbs}
	\begin{tabularx}{0.9\textwidth}{L l L}
		\toprule
		{\sc suffix} & {\sc meaning} & {\sc function} \\
		\midrule
		\ird{-ezak}    & until              & main clause holds until converb-clause event \\
		\ird{-eškady}  & when, around when  & loose temporal co-occurrence \\
		\ird{-eeby}    & since              & from the time of the converb-clause event \\
		\ird{-ešhoume} & as soon as         & immediately upon the converb-clause event \\
		\ird{-eebymy}  & after              & after the converb-clause event \\
		\ird{-esny}    & before             & before the converb-clause event \\
		\bottomrule
	\end{tabularx}
\end{table}

\section{Causal and concessive relations}

\subsection{Causal constructions}\label{sec:causal}


Causal subordination presents a reason or cause for the main-clause event
(`because X, Y'). The exponent of a dedicated causal converb is unspecified.

\subsection{Concessive constructions}\label{sec:concessive}

The concessive converb (\ird{-ouc}) introduces a conceded or contrasted clause
(`although X, Y'). It licenses the irrealis mood in the main clause when the
conceded content is contrary to expectations.

\section{Conditional constructions}\label{sec:conditional}

\subsection{The two conditional converbs}

Iridian distinguishes two conditional clause-linking relations: an open or
neutral conditional and an expected or predictive conditional. These are
expressed by two converb suffixes, summarised in Table~\ref{tab:conditional}.

\begin{table}[ht]
	\footnotesize
	\caption{Conditional converb suffixes.}\label{tab:conditional}
	\begin{tabularx}{\textwidth}{L l L L}
		\toprule
		{\sc suffix} & {\sc core meaning} & {\sc discourse profile} & {\sc notes} \\
		\midrule
		\ird{-ic}   & if             & open, neutral, contingent, generic &
			plain stem grade; contextual phonetics apply \\[4pt]
		\ird{-ezmy} & if / when      & expected, likely, scheduled, predictive &
			plain stem grade; see §\,\ref{sec:cond-nopala} \\
		\bottomrule
	\end{tabularx}
\end{table}

The conditional converb paradigm distinguishes open \ird{-ic} from expected
\ird{-ezmy}, without dedicated negative suffixes
(§\,\ref{sec:cond-negation}).

\subsection{Semantic distinction: open vs.\ expected conditional}
\label{sec:cond-semantics}

The opposition between \ird{-ic} and \ird{-ezmy} is grounded not in truth-value
alone but in the speaker's \emph{discourse stance toward fulfillment} of the
antecedent. The two forms can be distinguished as follows.

The suffix \ird{-ic} marks an \textbf{open conditional}: the speaker does not
commit to the fulfillment of the antecedent. It covers (a)~open real conditions,
where the outcome is uncertain; (b)~generic or habitual conditions,
where the antecedent expresses a recurrent situation type; (c)~hypothetical or
weakly counterfactual conditions, where speaker commitment is withheld; and
(d)~logical or instructional conditionals, where the if-clause functions as a
premise or rule rather than a prediction. In all these uses, the speaker
positions the antecedent as one possibility among others, without implying that
its occurrence is expected or imminent.

The suffix \ird{-ezmy} marks an \textbf{expected} or \textbf{predictive
conditional}: the speaker treats the fulfillment of the antecedent as likely,
scheduled, or inferable from context. It is used for (a)~events that are
anticipated or planned, (b)~natural-sequence conditionals in narrative or
procedural discourse where the antecedent is presented as a lead-in to the main
event, (c)~predictive \ird{if/when} constructions where occurrence is
presupposed or highly probable, and (d)~conditionals embedded in instructional
or procedural contexts where the antecedent is treated as a standard triggering
condition. The form is closer in discourse function to English \trsl{when} or
\trsl{once} in many of its uses, though the underlying relation remains
conditional.

The distinction thus pairs an epistemic dimension (speaker commitment to
antecedent fulfillment) with a discourse-pragmatic one (how the antecedent is
presented to the interlocutor). Neither suffix encodes temporal sequence
directly; temporal sequencing, where relevant, follows from the nature of the
main-clause event and its interaction with the mood system.

\pex \label{ex:cond-open}
\a
\begingl
\gla Marek zomenic, Janek pjáščev.//
\glb Marek leave-\Cond{} Janek eat-\Irr{}//
\glft \trsl{If Marek leaves, Janek might eat.} \lingcontext{open conditional; speaker uncommitted to Marek leaving; irrealis apodosis}//
\endgl
\a
\begingl
\gla Dana šegenic, novy zďíčev.//
\glb Dana ask-\Cond{} new bring-\Irr{}//
\glft \trsl{If Dana asks, [one] might bring a new one.} \lingcontext{generic/instructional; open conditional}//
\endgl
\xe

\pex \label{ex:cond-expected}
\a
\begingl
\gla Marek zomenezmy, Janek pjaštáž.//
\glb Marek leave-\ECond{} Janek eat-\Ctpl{}//
\glft \trsl{When/once Marek leaves, Janek will eat.} \lingcontext{expected conditional; departure treated as imminent or scheduled; indicative contemplative apodosis}//
\endgl
\a
\begingl
\gla Šaltezmy, dá pjasnou.//
\glb finish-\ECond{} \textsc{1sg} start-\Seq{}//
\glft \trsl{Once it finishes, I start.} \lingcontext{natural-sequence; procedural; antecedent is a standard triggering condition}//
\endgl
\xe

\subsection{Negation in conditional antecedents}\label{sec:cond-negation}

The conditional converb suffixes \ird{-ic} and \ird{-ezmy} have no dedicated
negative counterparts. Negation within an embedded finite indicative clause
follows the ordinary verbal pattern with \ird{za-}
(§\,\ref{sec:negation} of Chapter~\ref{ch:negation}), while the attachment and
scope of negation with a conditional converb itself are undocumented. The
finite negative mood suffixes \ird{-é}, \ird{-ku}, and \ird{-ir} do not belong
to the conditional converb paradigm. Conditionality itself does not license
negative indefinites: neither \ird{-ic} nor \ird{-ezmy} supplies the
negative operator required for concord (§\,\ref{sec:nm-negative-concord}).
The locality rule for finite clauses (§\,\ref{sec:complex-negative-concord})
does not establish concord across a converb boundary; the licensing of
negative indefinites inside non-finite conditional antecedents remains
unspecified.

\subsection{Morphophonology: no palatalisation with \ird{-ic}}\label{sec:cond-nopala}

The open conditional \ird{-ic} selects the plain consonant grade. Its initial
/ɪ/ does not by itself license the palatal/front or mixed alternants selected by
other suffixes (§\,\ref{sec:suffix-classes}). The suffix preserves plain
coronal stops, alveolar sibilants, and inherited velars; anticipatory
coarticulation before /ɪ/ does not constitute a categorical change in sibilant
place of articulation (§\,\ref{sec:palatal-sibilants}).

The expected conditional \ird{-ezmy} likewise selects the plain consonant
grade. Neither conditional suffix invokes the antipassive's replacement
pattern. Their grade membership is part of the morphological specification;
their historical development remains to be established independently.

\subsection{Clause type and pivot reference}\label{sec:cond-subj}

Conditional converb clauses are adjunct clauses within the general converb
system described in §\,\ref{sec:converb-chains}. Their internal alignment
follows the ordinary converb alignment: nominative-accusative case marking
applies inside the converb clause. An ordinary active has a direct-case actor
and indirect-case patient, while applicative and causative morphology select
the participant specified in §\,\ref{sec:converb-chains}. The
finite-verb-final rule holds: the conditional converb clause is non-final, and
the main clause supplies the sole finite verb.

Pivot reference in conditional converb clauses follows the default same-pivot
(SS) interpretation established for converbs generally. The absolutive pivot of
the \ird{-ic} or \ird{-ezmy} clause is by default coreferent with the absolutive
pivot of the matrix clause. Different-pivot (DS) readings require overt nominal
pivots or demonstratives in one or both clauses, or must be supported by
discourse context; Iridian has no grammaticalised switch-reference clitic (see
§\,\ref{sec:switch-reference}). In many everyday examples this will coincide
with same-agent reference because ordinary active converbs privilege the
actor; an applicative can instead maintain reference to its applied
participant. The generalisation remains one of pivot continuity.

\subsection{Interaction with irrealis in the main clause}\label{sec:cond-irrealis}

The choice between \ird{-ic} and \ird{-ezmy} correlates with, though does not
strictly determine, the mood of the apodosis. The open conditional \ird{-ic} is
the natural antecedent for main clauses in the irrealis (\ird{-ěv}) or the
negative irrealis (\ird{-ir}), because the speaker's non-commitment to
antecedent fulfillment is compatible with the irrealis framing of the consequent
as counterfactual, hypothetical, or weakly probable. This is the standard
configuration for hypothetical and counterfactual conditional constructions in
Iridian.

The expected conditional \ird{-ezmy} more naturally pairs with an indicative or
predictive main clause, because the antecedent is presented as likely or
scheduled and the consequent is therefore treated as a predictable outcome
rather than a mere possibility. Use of the irrealis with \ird{-ezmy} is marked
and implies a speaker retreat from the otherwise strong expectation that the
antecedent will be fulfilled; the construction then conveys something like
\trsl{even granting that X [which I expect], the result would merely be Y}.

For the irrealis system and its lexical negative counterpart \ird{-ir},
see §\,\ref{sec:mood} in Chapter~\ref{ch:verbs}.

\pex \label{ex:cond-irr}
\a
\begingl
\gla Marek zomenic, Janek pjáščev.//
\glb Marek leave-\Cond{} Janek eat-\Irr{}//
\glft \trsl{If Marek were to leave, Janek might eat.} \lingcontext{open conditional + irrealis apodosis; standard hypothetical}//
\endgl
\a
\begingl
\gla Marek zomenic, Janek pjaštir.//
\glb Marek leave-\Cond{} Janek eat-\NIrr{}//
\glft \trsl{If Marek were to leave, Janek would not eat.} \lingcontext{open conditional + negative irrealis apodosis}//
\endgl
\a
\begingl
\gla Marek zomenezmy, Janek pjaštáž.//
\glb Marek leave-\ECond{} Janek eat-\Ctpl{}//
\glft \trsl{When Marek leaves, Janek will eat.} \lingcontext{expected conditional + indicative contemplative apodosis; predicted outcome}//
\endgl
\xe

\subsection{Discourse function}\label{sec:cond-discourse}

The distinction between \ird{-ic} and \ird{-ezmy} is also sensitive to discourse
register and text type. The open conditional \ird{-ic} is the default form for
logical, generic, instructional, and contingent conditions. It appears
freely in expository prose, in legal and procedural texts where conditions are
stated as rules, and in any context where the speaker presents the antecedent as
one of several possible states of affairs.

The expected conditional \ird{-ezmy} is preferred in narrative, planning, and
predictive discourse, where the antecedent is treated as a natural lead-in to
the main event rather than as a genuine open alternative. It is
characteristically found in sequential narrative (\trsl{once X happens, then
Y}), in procedural instruction (\trsl{when the water boils, add the salt}), and
in predictive future statements where the speaker presents the antecedent as
part of an anticipated sequence of events. This makes \ird{-ezmy}
discourse-dependent in a way that \ird{-ic} is not: the choice of \ird{-ezmy} is
a stance-marking act that signals the speaker's commitment to the relevance and
imminence of the antecedent, whereas \ird{-ic} remains neutral on this
dimension.

\subsection{Formal analysis: converb features in Distributed Morphology}\label{sec:cond-dm}

In DM terms, the two conditional suffixes are Vocabulary Items inserted at the
defective converb head that underlies all adverbial converb constructions in
Iridian. This head differs from the ordinary finite verbal head in lacking
Asp.Align features (it does not drive perfective/imperfective interpretation or
the nominative-ergative case alternation), and it lacks the \synhead{MoodP}
layer that licenses directive, hortative, and irrealis readings. What the
conditional converb suffixes add is a set of clause-linking and discourse-stance
features beyond the bare \textsc{[conv]} feature shared by all adverbial
converbs. The two conditional exponents have the following feature specifications:

\pex
\a \ird{-ic}: \quad [\textsc{conv}, \textsc{cond}, \textsc{open}]
\hfill non-presuppositional; speaker uncommitted to antecedent fulfillment
\a \ird{-ezmy}: [\textsc{conv}, \textsc{cond}, \textsc{expected}]
\hfill presuppositionally live; antecedent treated as predictive or scheduled
\xe

The \textsc{open}/\textsc{expected} opposition encodes discourse stance
(§\,\ref{sec:cond-semantics}), while neither suffix marks negative polarity.
The expression of negation with a conditional converb is undocumented
(§\,\ref{sec:cond-negation}).

\section{Insubordinate conditional clauses}\label{sec:cond-insub}

A conditional converb clause ordinarily functions as an adverbial adjunct to a
matrix clause, as established in §\,\ref{sec:cond-subj}. Under specific
discourse conditions, however, a conditional converb clause may occur without an
overt apodosis. Such \textbf{insubordinate conditional clauses} are not analysed
here as ordinary adjuncts lacking their host through performance ellipsis.
Instead, the default analysis treats them as conditional clauses merged directly
into the discourse periphery as independent, licensed discourse units.

\subsection{Two structures}

The ordinary biclausal conditional follows the structure established in
§\,\ref{sec:converb-chains}: the conditional converb clause is merged as an
adjunct to the matrix \synhead{TP}/\synhead{IP}, scoping over the matrix
proposition (Figure~\ref{fig:cond-ordinary}).

\begin{figure}[ht]
\centering
\begin{forest}
for tree={align=center, l sep=10pt, s sep=14pt}
[\synhead{TP}
  [CondConvP
    [{$\ldots$}]
    [{V$_{\textsc{cond}}$}]
  ]
  [{T$'$}
    [{matrix material}]
    [{V$_{\textsc{finite}}$}]
  ]
]
\end{forest}
\caption{Ordinary conditional: converb clause as adjunct to \synhead{TP}.}
\label{fig:cond-ordinary}
\end{figure}

In the insubordinate use, no matrix \synhead{TP} is projected. Instead, the
conditional converb phrase is selected by a null framing head
(Frame\textsuperscript{0}) and the resulting FrameP is integrated into the
speech-act projection (Figure~\ref{fig:cond-insub}).

\begin{figure}[ht]
\centering
\begin{forest}
for tree={align=center, l sep=10pt, s sep=14pt}
[\synhead{SpeechActP}
  [FrameP
    [CondConvP
      [{$\ldots$}]
      [{V$_{\textsc{cond}}$}]
    ]
    [{Frame\textsuperscript{0}\\$\varnothing$}]
  ]
  [{SA\textsuperscript{0}}]
]
\end{forest}
\caption{Insubordinate conditional: converb phrase licensed at the discourse periphery.}
\label{fig:cond-insub}
\end{figure}

The null Frame\textsuperscript{0} head licenses the conditional converb phrase
as an independent discourse act without projecting an overt consequent. It
contributes a set of discourse instructions: adopt the antecedent as a live
hypothesis; evaluate the current or subsequent discourse under that hypothesis;
or supply the relevant consequence through pragmatic inference. In DM terms,
Frame\textsuperscript{0} is a null exponent subject to contextual licensing: it
selects a conditional converb phrase, requires a specific discourse environment
(warning, offer, topic-frame, or suspended inference), and is itself
interpretively inert---all the meaning comes from the converb clause it embeds
and the discourse context in which it is uttered.

\subsection{Pragmatic functions}

Insubordinate conditional clauses in Iridian cover several pragmatic functions.
In warning or caution use, the speaker presents a potential triggering condition
and leaves the consequence implicit, often to heighten its force. In offer or
invitation use, the antecedent names a condition available to the addressee,
with the result left for inference. In relevance-framing or topic-shift use, the
conditional clause restricts the discourse domain of what follows without
asserting an ordinary conditional proposition. In counterfactual suspension, the
speaker invites the interlocutor to consider a hypothetical space without
committing to any specific outcome.

\pex \label{ex:cond-insub}
\a
\begingl
\gla Zomenic\ldots//
\glb leave-\Cond{}//
\glft \trsl{If he leaves\ldots} \lingcontext{warning; consequence withheld; insubordinate}//
\endgl
\a
\begingl
\gla Dana šegenic\ldots//
\glb Dana ask-\Cond{}//
\glft \trsl{If Dana asks\ldots} \lingcontext{offer or relevance frame; insubordinate}//
\endgl
\a
\begingl
\gla Šaltezmy\ldots//
\glb finish-\ECond{}//
\glft \trsl{Once it finishes\ldots} \lingcontext{expected-conditional insubordinate; consequence strongly inferrable from context}//
\endgl
\xe

\subsection{Default form and the -ic/-ezmy contrast in insubordinate use}

The open conditional \ird{-ic} is the natural default form for insubordinate
conditionals. Because \ird{-ic} marks a non-presuppositional,
speaker-uncommitted condition ([\textsc{open}]), it is maximally compatible with
the suspended-inference pragmatics of insubordinate use: the speaker presents
the antecedent as a live possibility without foreclosing the consequent, and the
interlocutor is left to supply it. Warning, offer, and counterfactual-suspension
uses are therefore especially natural with \ird{-ic}.

The expected conditional \ird{-ezmy} can occur in standalone use, most naturally
in predictive or narrative contexts where the antecedent is treated as imminent
or scheduled (see (c) in~\ref{ex:cond-insub} above). This is the marked option:
because [\textsc{expected}] implies that the consequence is already accessible
in the discourse state, standalone \ird{-ezmy} tends to be restricted to cases
where the apodosis is strongly inferrable. Insubordinate \ird{-ezmy} can also
frame a topic or sequence in procedural discourse (\trsl{once it
finishes\ldots}, \trsl{when they arrive\ldots}), functioning as a
natural-sequence framing rather than a suspended warning.

\subsection{Secondary ellipsis analysis}

A secondary derivation is available where the consequence is locally recoverable
from the preceding discourse. In such cases, a conditional converb clause may be
followed by ellipsis of a recoverable matrix predicate at PF under discourse
identity:

\begin{figure}[ht]
\centering
\begin{forest}
for tree={align=center, l sep=10pt, s sep=14pt}
[\synhead{TP}
  [CondConvP
    [{$\ldots$}]
    [{V$_{\textsc{cond}}$}]
  ]
  [{T$'$\\{\footnotesize [deleted at PF]}}, roof]
]
\end{forest}
\caption{Secondary analysis: matrix clause deleted under recoverability.}
\label{fig:cond-ellipsis}
\end{figure}

This secondary derivation is appropriate where the apodosis is uniquely
recoverable from the immediate discourse context, and it should not be extended
as the default analysis. The primary reason is structural: conditional converbs
are already non-finite, reduced forms that project at most a
\synhead{vP}/\synhead{VoiceP}-level adjunct clause, not a full finite clause.
Positing a silent finite matrix after every insubordinate conditional would be a
stronger claim than the evidence requires, and would treat as syntactically
deficient what is in many discourse contexts a fully licensed, pragmatically
complete utterance. The insubordinate FrameP structure is therefore primary,
with the ellipsis derivation reserved for locally recoverable cases.

\subsection{Diagnostics for the insubordinate conditional}

Four properties distinguish insubordinate conditional clauses from
ordinary adjuncts whose matrix has been omitted in performance.

\textit{Prosodic independence.} An insubordinate conditional carries its own
intonation contour and closes an intonational phrase in spoken Iridian. An
ordinary dependent converb can share an intonational phrase with its matrix,
although internal phonological phrasing remains available
(§\,\ref{sec:converb-phrasing}). Prosodic independence does not require
silence at the boundary or prescribe a particular terminal melody.

\textit{Independent illocutionary force.} An insubordinate conditional can carry
its own illocutionary force---warning, challenge, offer--- independently of any
overt apodosis. A dependent converb adjunct has no independent force; it merely
conditions the illocutionary force of the matrix.

\textit{Non-recovery of a unique apodosis.} Some insubordinate conditionals do
not recover a unique apodosis at all. A warning framing (\trsl{if he
comes\ldots}) or a counterfactual suspension (\trsl{if only\ldots}) opens an
inferential space rather than pointing to a specific recoverable proposition. A
matrix-ellipsis analysis is too strong in such cases.

\textit{Preference for \ird{-ic}.} Insubordinate conditionals show a strong
preference for the open conditional suffix over \ird{-ezmy}. This follows
directly from the feature analysis: [\textsc{open}] conditional stance is
maximally compatible with suspended-inference pragmatics; [\textsc{expected}]
stance implies that the consequence is already accessible in the discourse
state, restricting standalone use to highly inferential contexts.

\section{Purposive constructions}\label{sec:purposive}

The purposive converb (\ird{-it}) expresses purpose (`in order to X') and serves
as the complement of modal, desiderative, and motion predicates. Its consonant
alternants are specified by stem class rather than generated by a general
palatal-grade rule: \ird{pjašcit}, for example, is the listed Class~I
purposive alternant of \ird{pjašt-}. This /ʃt/--/ʃts/ correspondence does not
extend to every stem before \ird{-it}. As the canonical non-finite
event-complement form (§\,\ref{sec:th-compl-distinction}), the purposive converb
is licensed wherever the matrix predicate selects an event regarded as
projected, volitional, incipient, or purpose-oriented rather than a full
proposition.

\subsection{The intentional future reanalysis}\label{sec:purp-intentional}

In colloquial speech, the purposive converb has been reanalysed as a
future/intentional marker when no governing verb is present. The first-person
subject pronoun \ird{dá} is typically retained overtly in this construction, as
the purposive suffix alone does not encode person:

\ex
\begingl
\gla Dá travě pjašcit.//
\glb \textsc{1sg} bread.\Ind{} eat-\Purp{}//
\glft \trsl{I'm gonna eat bread.} \lingcontext{colloquial intentional future}//
\endgl
\xe

\section{The purposive-motion construction}\label{sec:purp-motion}

A productive extension of the event-complement pattern is the
\term{purposive-motion construction}, in which a motion verb---\ird{stoj-} `go',
\ird{prit-} `come', \ird{žedet-} `go out', or a closely related verb of directed
or undirected motion---serves as the matrix predicate and the purposive converb
supplies the complement naming the purpose of the motion. The construction
expresses motion undertaken in order to bring about a stated event:

\begin{center}
\term{[\textsc{purp.conv} complement] \quad [motion verb\textsubscript{fin}]}
\end{center}

\noindent No new structural mechanism is required. The purposive converb already
functions as the complement of modal, desiderative, and inchoative matrix
predicates; motion predicates constitute a natural extension of this
complement-selecting class (§\,\ref{sec:th-compl-distinction}). The
complement-final order follows Iridian's general head-final architecture, and
same-pivot continuity across the converb boundary follows the rule established
in §\,\ref{sec:switch-reference}.

Three overlapping semantic domains are distinguished in the subsections that
follow.

\subsection{Literal purposive motion}\label{sec:purp-motion-literal}

In its core use the construction describes concrete, physically real motion
undertaken for a stated purpose. The mover departs or arrives; the purposive
converb names the reason.

\ex\label{ex:purp-swim}%
\begingl
\gla Tězně knářit stójek.//
\glb lake.\Ind{} swim-\Purp{} go-\Pf{}//
\glft \trsl{(They) went swimming in the lake.}//
\endgl
\xe

\ex\label{ex:purp-meet}%
\begingl
\gla Josifam všetit žedétek.//
\glb Josif.\Ind{} meet-\Purp{} go.out-\Pf{}//
\glft \trsl{(He) went out to meet Josif.}//
\endgl
\xe

\ex\label{ex:purp-smoke}%
\begingl
\gla Raznit žedétek.//
\glb smoke-\Purp{} go.out-\Pf{}//
\glft \trsl{(He) went out to smoke.}//
\endgl
\xe

In each example the matrix mover is null and recovered from discourse;
the implicit mover of the purposive clause shares its reference through
the control relation. Discourse recoverability and pivot control are
therefore separate components of the interpretation. Case marking inside the purposive
complement observes nominative-accusative alignment
(§\,\ref{sec:converb-chains}): the sole argument of an intransitive converb such
as \ird{knář-} `swim' is the direct-case absolutive, and the patient of a
transitive converb such as \ird{všet-} `meet' appears in the indirect (oblique)
case---hence \ird{Josifam} (Josif.\Ind{}) in (\ref{ex:purp-meet}), and
\ird{tězně} (lake.\Ind{}) as the locative argument of `swim' in
(\ref{ex:purp-swim}). The matrix motion verbs \ird{stoj-} and \ird{žedet-} are
themselves intransitive; the mover is their sole participant and surfaces in the
absolutive (direct) case when overt. The ordinary intransitive perfective suffix \ird{-ek} is
inserted at the clause-final finite position, concentrating aspect and alignment
at the right edge of the clause chain in accordance with the finite-verb-final
rule (§\,\ref{sec:converb-chains}).

The stem \ird{razeň-} `smoke' selects reduced \ird{razn-} before the
purposive suffix, yielding \ird{raznit} (§\,\ref{sec:purp-motion-structure}).

\subsection{Prospective motion and purposive
foregrounding}\label{sec:purp-motion-prospective}

A closely related reading arises when motion is directed toward the speech-act
participants and the purposive content is communicatively prominent. The
physical motion recedes and the purpose is foregrounded:

\ex\label{ex:purp-come}%
\begingl
\gla Že pomadit by prítek?//
\glb \textsc{1sg}.\Ind{} talk.to-\Purp{} \textsc{q} come-\Pf{}//
\glft \trsl{Did you come here to talk to me?}//
\endgl
\xe

In (\ref{ex:purp-come}), \ird{prit-} `come' encodes deictic arrival toward the
speaker. The polar question particle \ird{by} realises matrix Pol, while
the selected purposive complement \ird{že pomadit} `to talk to me' undergoes
ordinary placement as a whole before the matrix particle field. Its internal
argument is not independently extracted into the matrix clause; \ird{že} is the indirect-case (oblique) form of the first-person singular
pronoun, appearing here as the patient of the transitive converb \ird{pomad-}.
The null subject (\textsc{2sg}) is recoverable from the interrogative context.
The reading is licensed directly by the semantics of event complementation: the
purposive converb already presents its complement as a prospective,
purpose-oriented event rather than a proposition
(§\,\ref{sec:th-compl-distinction}), and no reanalysis is required to produce it
with deictic-motion verbs.

\subsection{Colloquial extension: intention and
imminence}\label{sec:purp-motion-colloquial}

In informal and colloquial registers, \ird{stoj-} and \ird{prit-} may appear
with reduced or bleached motion semantics, and the purposive-motion construction
approximates intention or near-future imminence. Where the motion component
contributes little to the interpretation, the purposive converb can carry the
prospective content without an overt matrix verb
(§\,\ref{sec:purp-intentional}). This register-marked construction contrasts
with the concrete motion reading, in which departure or arrival remains part of
the asserted event.

\subsection{Structural analysis}\label{sec:purp-motion-structure}

Within the Distributed Morphology framework adopted throughout this grammar, the
motion verbs \ird{stoj-}, \ird{prit-}, \ird{žedet-}, and \ird{nočk-} are
ordinary Vocabulary Items inserted at V\textsuperscript{0}. They are
complement-selecting lexical predicates, not auxiliaries; they are not merged at
Asp\textsuperscript{0} or any higher functional head. The purposive converb
suffix \ird{-it} is the exponent of a converb head (Conv\textsuperscript{0})
bearing [\textsc{+purposive}] features; it heads the purposive complement phrase
merged in the complement position of the motion verb's VP projection. The
resulting structure is identical in all relevant respects to any other
event-complement construction in Iridian: the motion verb selects a non-finite
purposive complement in exactly the way that desiderative or inchoative
predicates do, and no special licensing rule is required.

Morphophonology within the construction follows independently motivated rules.
For the fleeting-vowel stem \ird{razeň-} `smoke', the purposive selects
\ird{razn-}; its hard nasal is part of that lexical bound grade. Appending
\ird{-it} gives \ird{raznit}, with short unstressed \ird{e} absent and the
stressed vowel preserved. The clause-final intransitive perfective
suffix \ird{-ek} is inserted at
Asp\textsuperscript{0}\,$\circ$\,Align\textsuperscript{0} by the general rules
of the aspect-alignment paradigm, with intransitive motion verbs receiving
absolutive alignment (direct-case mover). No additional morphological licensing
is specific to the purposive-motion construction.

\subsection{Limits of the construction: contrast with aspectual
periphrasis}\label{sec:purp-motion-limits}

The construction might invite comparison with Japanese \textit{V-te iku} /
\textit{V-te kuru}, in which motion verbs have grammaticalized into markers of
directional aspect encoding result, continuation, or change of state through
time. Such a development is not available in Iridian for principled
architectural reasons.

Iridian concentrates viewpoint and alignment in the clause-final finite verb.
The aspectual inventory---continuous, progressive, perfective, retrospective,
contemplative---is encoded by the
Asp\textsuperscript{0}\,$\circ$\,Align\textsuperscript{0} complex at the right
edge of the clause chain. A grammaticalized \textit{converb + motion verb}
aspectual periphrasis would duplicate this work, competing with the existing
system rather than harmonizing with it. The language therefore has neither the
structural need nor the available niche for a general directional-aspectual
auxiliary system of the Japanese type.

The purposive-motion construction is accordingly kept narrow. \ird{Stoj-} and
\ird{prit-} do not develop the broad continuative-change semantics of Japanese
\textit{iku} and \textit{kuru}, and remain complement-selecting lexical
predicates at every level of the synchronic grammar. Any auxiliary-like flavour
in colloquial speech is a matter of usage and ongoing grammaticalisation
(§\,\ref{sec:purp-motion-colloquial}), not of the core syntax.

A final example with an undirected motion verb confirms that the construction is
not restricted to displacement predicates:

\ex\label{ex:purp-run}%
\begingl
\gla Mé autobusam zrolit nóčkek.//
\glb \textsc{1pl} bus.\Ind{} catch.up-\Purp{} run-\Pf{}//
\glft \trsl{We ran to catch up with the bus.}//
\endgl
\xe

Here \ird{nočk-} `run' is the matrix motion verb; it denotes undirected
locomotion rather than displacement, yet it selects a purposive complement
without difficulty. The object of the converb \ird{zrol-} `catch up with' is
\ird{autobus} `bus', whhich takes the indirect case by virtue of its patient
role inside the nominative-accusative converb clause. The first-person plural
subject \ird{mé} is overtly expressed, as nothing in the preceding discourse
supplies a recoverable plural topic.

\section{Contrastive linkage and coordination}

\subsection{Contrastive constructions}\label{sec:contrastive}

The contrastive converb (\ird{-ema}) introduces a contrasting clause (`X,
but/whereas Y'), covering simple contrast, adversative contrast, and parallel
contrast.

\subsection{Coordination}\label{sec:coordination}


Coordination is expressed by \ird{a} (additive), \ird{še} (additive, also
complementizer), \ird{má} / \ird{ozná} (adversative), and a set of disjunctives.


Conjunctions link coordinated constituents and clauses
(§\,\ref{sec:coordination}), whereas discourse connectives relate an utterance
to the preceding exchange. The connective \ird{nu} is treated in
§\,\ref{sec:spoken-nu}; clause-final discourse particles, which stand outside
the preverbal sequence, are discussed in §\,\ref{sec:disc-final-particles}.


\section{State-frame adverbial clauses}
\label{sec:state-frame}

Iridian possesses a special non-finite suffix -eň, used with nominal predicates
and with a restricted class of stative predicates to form state-frame adverbial
clauses. The construction expresses meanings such as ‘while X’, ‘when in the
state of X’, and ‘as X’. It is functionally parallel to the ordinary converb
system in that it forms a preverbal adverbial dependent clause, but it is not
part of the regular converb inventory and does not replace an ordinary Asp.Align
morpheme in the verbal paradigm.

Semantically, -eň backgrounds a classification or state as the temporal or
circumstantial frame for the matrix event. The clause does not assert entry into
the state, progression toward it, or completion of a transition. In this it
contrasts sharply with neh-, whose function is explicitly transitional and
inchoative: neh- introduces an event of becoming and projects a full verbal
spine under \synhead{AspP}. A form such as \ird{byl nesu} would therefore mean
something like ‘while becoming a child’ or ‘while entering child-status’, not
the intended simple background reading ‘while a child’. The state-frame converb
-eň is the dedicated exponent of the latter meaning.

Syntactically, a -eň clause behaves like other Iridian adverbial subordinate
clauses. It precedes the matrix clause in accordance with general head-finality,
and its interpretation normally follows the same-pivot continuity pattern unless
context or overt nominal material indicates otherwise. Because the construction
is not an ordinary finite clause, it bears no independent aspect or alignment
morphology; the matrix finite verb remains the sole carrier of aspect, mood, and
alignment in the chain, in accordance with the general finite-verb-final rule
for Iridian subordination.


The suffix is primarily used with class-membership nouns and age/status
predicates: byl-eň ‘while a child, as a child’ kral-eň ‘while king, as king’
vdov-eň ‘while a widow’

The phasal particle \ird{so} may occur with the state-frame converb -eň. In such cases
it scopes over the predicative state expressed by the converb and marks that
state as still holding at the reference time of the subordinate clause. Thus
byleň means ‘while a child’, whereas \ird{so byleň} means ‘while still a child’. Any
implication that the state no longer holds at speech time arises pragmatically
from discourse perspective, not from the particle itself.
